\section{Conclusion}

Ce projet illustre le passage d'un prototype GraphRAG fonctionnel à un système dont la fiabilité a été construite incident par incident, chacun découvert par un test réel et corrigé à sa racine plutôt que contourné. L'architecture retenue -- routage d'intention, planification d'outils, récupération hybride Neo4j/Qdrant, fusion pondérée par des règles métier explicites, reclassement sémantique, filtre de pertinence adapté à chaque type de preuve disponible, garde-fou de sortie, et désormais mémoire conversationnelle par session -- traduit directement les trois exigences fixées au départ~: exactitude vérifiable, bilinguisme réel, et fonctionnement entièrement local sur une infrastructure CPU maîtrisée.

Les corrections apportées ne relèvent pas d'ajustements cosmétiques~: un filtre de pertinence structurellement incapable de rejeter une question hors domaine, une couverture de données nulle sur deux catégories de population entières, ou une réponse de repli qui omettait le contenu réel des foires aux questions et des programmes, sont autant de défauts qui, non corrigés, auraient directement compromis la confiance que l'exigence d'exactitude cherchait précisément à garantir. Leur résolution méthodique, documentée dans ce rapport avec le symptôme, le diagnostic, la solution et la vérification propres à chacun, constitue la meilleure garantie apportée à ce jour de la fiabilité du système.

Les limitations qui subsistent -- gestion de la négation, latence sur ce CPU partagé, absence de superviseur de processus et de gestion de versions -- sont explicitement documentées plutôt que dissimulées, et constituent la feuille de route naturelle de la prochaine itération du projet.
